\section{Optimización del clasificador}

\subsection{Baseline de referencia}

EfficientNet-B0 se mantuvo como punto de partida porque había obtenido el mejor
macro-F1 de la Etapa 1. En la corrida nueva no se volvió a entrenar toda la red:
se retiró su clasificador, se congelaron los 4\,007\,548 parámetros del
backbone ImageNet y se extrajo un vector de 1\,280 características por imagen.
Sobre esos vectores se ajustó una regresión logística multiclase incremental
con SGD.

El baseline interno de esta etapa usa $\alpha=10^{-4}$, tasa constante 0.01,
penalización L2, batch 256, promedio de pesos, normalización L2 de embeddings y
20 épocas. No debe confundirse con el EfficientNet-B0 end-to-end histórico. El
primero sirve para aislar el efecto del tuning bajo el nuevo protocolo; el
segundo conserva la referencia de 0.9146 macro-F1 sobre el corpus y test de la
Etapa 1.

\subsection{Espacio de hiperparámetros}

Optuna se eligió porque el muestreador TPE usa los trials previos para concentrar
la búsqueda y porque el pruning puede detener combinaciones que ya muestran un
macro-F1 poco competitivo \citep{akiba2019optuna}. Se ejecutaron 30 trials con
semilla 42. El objetivo fue macro-F1 en las 5\,015 imágenes de validación; el
holdout no participó.

La búsqueda cubrió $\alpha$ entre $10^{-6}$ y $10^{-3}$ en escala logarítmica,
tasa inicial entre $3\times10^{-4}$ y $5\times10^{-2}$ también logarítmica,
penalización L2 o elastic net, proporción L1 entre 0 y 0.5 cuando correspondía,
batch de 128, 256 o 512, promedio de pesos activado o no, embeddings con o sin
normalización L2 y potencia de balanceo 0, 0.5 o 1. Esta última transforma los
pesos inversos de clase: cero no rebalancea, 0.5 aplica raíz cuadrada y uno usa
la inversa completa.

\begin{table}[H]
  \centering
  \caption{Configuración del baseline interno y combinación seleccionada por
  Optuna.}
  \label{tab:tuning-config}
  \small
  \begin{tabular}{lcc}
\toprule
\textbf{Hiperparámetro} & \textbf{Baseline} & \textbf{Mejor valor} \\
\midrule
alpha & 0.0001 & 4.2836449125618075e-06 \\
eta0 & 0.01 & 0.0006059127453715242 \\
penalty & l2 & l2 \\
l1\_ratio & 0.15 & 0.0 \\
batch\_size & 256 & 512 \\
balance\_power & 0.0 & 0.5 \\
average & True & False \\
feature\_norm & l2 & none \\
epochs & 20 & 20 \\
\bottomrule
\end{tabular}

\end{table}

El pruner fue \emph{MedianPruner}, con cinco trials iniciales y cuatro épocas de
calentamiento. Cada época reportó macro-F1, de modo que una configuración podía
terminar antes de las veinte épocas. La tabla de trials, el historial por época
y la importancia de parámetros se guardaron directamente desde el estudio.

\subsection{Resultados del tuning}

\begin{figure}[H]
  \centering
  \includegraphics[width=0.88\textwidth]{optuna_history.pdf}
  \caption{Macro-F1 de los trials completos y mejor valor acumulado. La línea
  discontinua corresponde al baseline de la cabeza, no al modelo histórico.}
  \label{fig:optuna-history}
\end{figure}

\begin{figure}[H]
  \centering
  \includegraphics[width=0.86\textwidth]{optuna_importance.pdf}
  \caption{Importancia relativa calculada por Optuna dentro del espacio
  explorado. Una importancia baja no prueba que el hiperparámetro sea inútil
  fuera de esos rangos.}
  \label{fig:optuna-importance}
\end{figure}

\begin{table}[H]
  \centering
  \caption{Cinco mejores trials completos según macro-F1 de validación.}
  \label{tab:top-trials}
  \begin{tabular}{rrl}
\toprule
\textbf{Trial} & \textbf{Macro-F1 val.} & \textbf{Estado} \\
\midrule
26 & 0.8887 & COMPLETE \\
11 & 0.8877 & COMPLETE \\
12 & 0.8877 & COMPLETE \\
23 & 0.8877 & COMPLETE \\
22 & 0.8865 & COMPLETE \\
\bottomrule
\end{tabular}

\end{table}

\begin{table}[H]
  \centering
  \caption{Efecto observado del tuning sobre la misma división y backbone.}
  \label{tab:baseline-tuned}
  \begin{tabular}{lrrr}
\toprule
\textbf{Configuración} & \textbf{Accuracy} & \textbf{Macro-F1} & \textbf{F1 ponderado} \\
\midrule
baseline probe & 0.8742 & 0.7799 & 0.8685 \\
optuna tuned probe & 0.9382 & 0.8887 & 0.9383 \\
\bottomrule
\end{tabular}

\end{table}

Optuna completó 17 trials y podó 13. El mejor fue el trial 26 en la
numeración mostrada al lector (índice interno 25): penalización L2,
$\alpha=4.28\times10^{-6}$, tasa 0.000606, batch 512, potencia de balanceo
0.5, sin promedio de pesos y sin normalización L2 de embeddings. El patrón de
los mejores trials fue consistente: una regularización pequeña, tasa menor que
la del baseline y balanceo por raíz cuadrada.

El macro-F1 subió de 0.7799 a 0.8887, una ganancia absoluta de 0.1088 sobre la
misma división. Accuracy pasó de 0.8742 a 0.9382 y el cambio mayor entre los
componentes de macro-F1 apareció en recall macro: 0.7516 a 0.8914. Precision
macro también aumentó, de 0.8511 a 0.8864. El baseline no fallaba solamente por
falta de separación; estaba favoreciendo demasiado a las clases grandes. La
raíz de los pesos inversos corrigió parte de ese sesgo sin llegar al balanceo
completo, que varios trials llevaron a resultados peores.

La Figura~\ref{fig:optuna-importance} atribuye la mayor importancia a la
penalización (0.261), seguida por tasa de aprendizaje (0.189), promedio de
pesos (0.186) y potencia de balanceo (0.176). Alpha quedó en 0.038 dentro de
esta búsqueda. Esto no significa que la regularización sea irrelevante: los
trials competitivos ya se concentraron en valores pequeños. La lectura útil es
que no hizo falta aumentar la capacidad del backbone. La ganancia vino de cómo
la cabeza convirtió sus características en fronteras para clases desiguales.
